% ===== PRÉAMBULE CANONIQUE — à coller VERBATIM en tête de CHAQUE .tex =====
\documentclass[11pt,a4paper]{article}
\usepackage[utf8]{inputenc}\usepackage[T1]{fontenc}\usepackage{lmodern}
\usepackage[french]{babel}
\usepackage{amsmath,amssymb,mathtools}\usepackage{icomma}
\usepackage[version=4]{mhchem}          % \ce{...} : équations & formules
\usepackage{siunitx}\sisetup{locale=FR} % \qty{}{}, \si{}, \ang{}
\usepackage{chemfig}                    % \chemfig{...} : structures développées
\usepackage{xcolor}\usepackage{colortbl}
\usepackage{tikz}\usepackage{pgfplots}\pgfplotsset{compat=1.18}
\usepackage[most]{tcolorbox}\usepackage{enumitem}\usepackage{array,booktabs}
\usepackage[a4paper,margin=2.2cm]{geometry}
\usetikzlibrary{arrows.meta,calc,angles,quotes,positioning,patterns,backgrounds,shapes.geometric}
\definecolor{primary}{RGB}{222,49,99}\definecolor{primarydark}{RGB}{150,22,55}
\definecolor{defcol}{RGB}{0,114,178}\definecolor{methcol}{RGB}{52,92,148}
\definecolor{excol}{RGB}{0,135,90}\definecolor{attncol}{RGB}{200,40,55}
\tcbset{boxbase/.style={enhanced,breakable,boxrule=0.9pt,arc=2.5pt,
  left=3mm,right=3mm,top=2.3mm,bottom=2.3mm,fonttitle=\bfseries,coltitle=white}}
% --- boîtes TUTORIEL (noms EXACTS, ne pas renommer) ---
\newtcolorbox{cle}[1][]{boxbase,colback=defcol!6,colframe=defcol,
  title={\textsc{À retenir}\ifx\relax#1\relax\relax\else\textmd{ : \itshape #1}\fi}}
\newtcolorbox{methode}[1][]{boxbase,colback=methcol!7,colframe=methcol,sharp corners,
  title={\textsc{Méthode}\ifx\relax#1\relax\relax\else\textmd{ --- \itshape #1}\fi}}
\newtcolorbox{exemplebox}[1][]{boxbase,colback=excol!5,colframe=excol,
  title={\textsc{Exemple}\ifx\relax#1\relax\relax\else\textmd{ : \itshape #1}\fi}}
\newtcolorbox{piege}[1][]{boxbase,colback=attncol!6,colframe=attncol,
  title={\textsc{Piège}\ifx\relax#1\relax\relax\else\textmd{ : \itshape #1}\fi}}
\newtcolorbox{remarque}[1][]{boxbase,colback=black!4,colframe=black!45,
  title={\textsc{Remarque}\ifx\relax#1\relax\relax\else\textmd{ : \itshape #1}\fi}}
% --- boîtes EXERCICE / CORRIGÉ (noms EXACTS, auto-comptées) ---
\newtcolorbox[auto counter]{exo}[1][]{boxbase,colback=primary!4,colframe=primary,
  title={\textsc{Exercice}~\thetcbcounter\ifx\relax#1\relax\relax\else\textmd{ --- \itshape #1}\fi}}
\newtcolorbox[auto counter]{corrige}[1][]{boxbase,colback=excol!4,colframe=excol,
  title={\textsc{Corrigé}~\thetcbcounter\ifx\relax#1\relax\relax\else\textmd{ --- \itshape #1}\fi}}
\newcommand{\R}{\mathbb{R}}\newcommand{\N}{\mathbb{N}}\newcommand{\Z}{\mathbb{Z}}\newcommand{\ds}{\displaystyle}
% =========================================================================
\begin{document}
% THEME_TITLE: Relation entre Kps et solubilité
\sisetup{per-mode=symbol}
\section*{\color{primarydark}Thème 3 — Relation entre $K_{\mathrm{ps}}$ et solubilité $s$}

$K_{\mathrm{ps}}$ et $s$ décrivent le \emph{même} état de saturation : $s$ est un nombre de moles
dissoutes, $K_{\mathrm{ps}}$ un produit de concentrations. Ce thème apprend à passer de l'un à l'autre
— la mécanique la plus « calculatoire » de la fiche.

\begin{methode}[exprimer $K_{\mathrm{ps}}$ en fonction de $s$]
\begin{enumerate}[leftmargin=1.7em,topsep=2pt,itemsep=1.5pt]
  \item Écrire l'équation de dissolution ($p$ cations, $q$ anions).
  \item À l'équilibre : $[\ce{M^{q+}}]=p\,s$ et $[\ce{X^{p-}}]=q\,s$.
  \item Reporter dans $K_{\mathrm{ps}}=[\ce{M^{q+}}]^{p}[\ce{X^{p-}}]^{q}$ et simplifier :
    \[ \boxed{\,K_{\mathrm{ps}} = (p\,s)^{p}(q\,s)^{q} = p^{\,p}\,q^{\,q}\;s^{\,p+q}\,} \]
\end{enumerate}
\end{methode}

\begin{cle}[le tableau à connaître par cœur]
\renewcommand{\arraystretch}{1.5}
\begin{center}\small
\begin{tabular}{c c c c}
\toprule
\rowcolor{defcol!12}
\textbf{Type} & \textbf{Exemples} & \textbf{$K_{\mathrm{ps}}$ en fonction de $s$} & \textbf{$s$ en fonction de $K_{\mathrm{ps}}$}\\
\midrule
$1{:}1$ & \ce{AgCl}, \ce{CaCO3} & $K_{\mathrm{ps}}=s^{2}$ & $s=\sqrt{K_{\mathrm{ps}}}$\\
$1{:}2$ ou $2{:}1$ & \ce{PbBr2}, \ce{Ag2CrO4} & $K_{\mathrm{ps}}=4\,s^{3}$ & $s=\sqrt[3]{K_{\mathrm{ps}}/4}$\\
$1{:}3$ ou $3{:}1$ & \ce{Fe(OH)3}, \ce{Ag3PO4} & $K_{\mathrm{ps}}=27\,s^{4}$ & $s=\sqrt[4]{K_{\mathrm{ps}}/27}$\\
$3{:}2$ & \ce{Ca3(PO4)2} & $K_{\mathrm{ps}}=108\,s^{5}$ & $s=\sqrt[5]{K_{\mathrm{ps}}/108}$\\
\bottomrule
\end{tabular}
\end{center}
\end{cle}

\begin{exemplebox}[deux calculs types]
\textbf{\ce{Fe(OH)3}} (type $1{:}3$) : $[\ce{Fe^{3+}}]=s$, $[\ce{OH-}]=3s$, donc
$K_{\mathrm{ps}}=(s)(3s)^{3}=27\,s^{4}$.

\smallskip
\textbf{\ce{AgCl}} (type $1{:}1$) : $K_{\mathrm{ps}}=s^{2}$, donc
$s=\sqrt{K_{\mathrm{ps}}}=\sqrt{\num{1.8e-10}}=\qty{1.34e-5}{\mole\per\litre}$.
\end{exemplebox}

\begin{piege}[ne jamais comparer directement des $K_{\mathrm{ps}}$ de types différents !]
Entre sels de \textbf{même type}, le plus petit $K_{\mathrm{ps}}$ est le moins soluble. Mais entre
\textbf{types différents}, il faut repasser par $s$ :
\[ K_{\mathrm{ps}}(\ce{AgCl})=\num{1.8e-10}\ (1{:}1)\Rightarrow s=\qty{1.34e-5}{\mole\per\litre}, \]
\[ K_{\mathrm{ps}}(\ce{Ag2CrO4})=\num{1.1e-12}\ (2{:}1)\Rightarrow s=\qty{6.5e-5}{\mole\per\litre}. \]
\ce{Ag2CrO4} a le $K_{\mathrm{ps}}$ le \emph{plus petit} mais est \emph{plus} soluble que \ce{AgCl} !
\end{piege}

\begin{piege}[{$s\neq[\text{ion}]$ pour les sels non $1{:}1$}]
La solubilité $s$ n'est \textbf{pas} égale à la concentration d'un ion, sauf pour un ion de coefficient
$1$. Pour \ce{PbBr2} : $s=[\ce{Pb^{2+}}]$ mais $[\ce{Br-}]=2s$. On raisonne toujours sur $s$, puis on
multiplie par le coefficient.
\end{piege}

\begin{remarque}[à retenir sur les racines]
Pour un sel $1{:}1$ le calcul de $s$ est une simple \emph{racine carrée} ; dès qu'un coefficient dépasse
$1$, c'est une racine cubique, quatrième\dots\ On laisse souvent le résultat sous forme de racine
(par exemple $s=\sqrt[3]{K_{\mathrm{ps}}/4}$), comme au concours.
\end{remarque}
\end{document}
